\section{Bab I: Penerapan Nilai-Nilai Pancasila}
\subsection{A. Asal Mula dan Kedudukan Pancasila}
Pancasila adalah pandangan hidup bangsa yang nilai-nilainya telah hidup dalam masyarakat Indonesia sejak sebelum merdeka. Semboyan \textit{Bhinneka Tunggal Ika} diambil dari Kitab Sutasoma karya Mpu Tantular pada masa Kerajaan Majapahit.

Menurut Kaelan (2020), secara kausalitas asal mula Pancasila dibagi menjadi empat:
\begin{enumerate}
    \item \textbf{Kausa Materialis (Asal Mula Bahan):} Unsur-unsur Pancasila digali dari nilai adat istiadat, kebudayaan, serta nilai religius bangsa Indonesia.
    \item \textbf{Kausa Formalis (Asal Mula Bentuk):} Pancasila dirumuskan dan dibahas oleh BPUPK dan disahkan oleh PPKI.
    \item \textbf{Kausa Efisien (Asal Mula Karya):} Tindakan PPKI sebagai pembentuk negara yang mengesahkan Pancasila menjadi dasar negara yang sah.
    \item \textbf{Kausa Finalis (Asal Mula Tujuan):} Pancasila dirumuskan oleh pendiri negara dengan tujuan untuk ditetapkan sebagai dasar negara.
\end{enumerate}

Terdapat tiga tataran nilai dalam ideologi Pancasila, yaitu: \textbf{nilai dasar} (fundamental dan universal), \textbf{nilai instrumental}, dan \textbf{nilai praksis}.

\subsection{B. Pembudayaan Pancasila dan Investasi Karakter}
Pembudayaan Pancasila bertujuan untuk mewujudkan tujuan negara dengan mengedepankan potensi kecerdasan (\textit{civic intelligence}). Ir. Sukarno menggagas \textbf{Revolusi Mental} (\textit{Nation and Character Building}) untuk membangun mentalitas bangsa. 

Dalam pendidikan karakter, ajaran keteladanan Ki Hajar Dewantara sangat penting diterapkan:
\begin{itemize}
    \item \textit{Ing Ngarso Sung Tulodo}: Di depan memberi teladan.
    \item \textit{Ing Madyo Mangun Karso}: Di tengah membangun kehendak/semangat.
    \item \textit{Tut Wuri Handayani}: Di belakang memberikan dorongan moral.
\end{itemize}

\subsection{C. Strategi Perubahan dan Pengamalan}
Menurut Yudi Latif (2020), aktualisasi nilai Pancasila dapat ditempuh melalui lima jalur utama: (1) Jalur pemahaman, (2) Jalur kerukunan, (3) Jalur keadilan, (4) Jalur pelembagaan, dan (5) Jalur keteladanan.

Contoh pengamalan Pancasila dalam kehidupan sehari-hari:
\begin{itemize}
    \item \textbf{Lingkungan Keluarga:} Menghormati keluarga, menyelesaikan masalah dengan diskusi, dan saling membantu.
    \item \textbf{Lingkungan Sekolah:} Mematuhi tata tertib, toleransi antarsiswa, dan mencegah perundungan (\textit{bullying}).
    \item \textbf{Lingkungan Masyarakat:} Rukun dengan tetangga, mengutamakan musyawarah, dan tolong-menolong.
    \item \textbf{Lingkungan Bangsa dan Negara:} Toleransi keberagaman, mematuhi hukum yang berlaku, dan menyampaikan aspirasi secara bertanggung jawab.
\end{itemize}